%
% zahlengerade.tex -- template for standalon tikz images
%
% (c) 2021 Prof Dr Andreas Müller, OST Ostschweizer Fachhochschule
%
\documentclass[tikz]{standalone}
\usepackage{amsmath}
\usepackage{times}
\usepackage{txfonts}
\usepackage{pgfplots}
\usepackage{csvsimple}
\usetikzlibrary{arrows,intersections,math}
\begin{document}
\def\skala{0.7}
\begin{tikzpicture}[>=latex,thick,scale=\skala]

% Zahlengerade von -4 bis 10 mit Pfeil
\draw[->] (-4.7, 0) -- (10.5, 0);
% Teilstriche fuer alle ganzen Zahlen von -4 bis 10
\foreach \x in {-4,-3,...,10} {
    \draw (\x, 0.08) -- (\x, -0.08);
    \node[above=3pt, font=\scriptsize] at (\x, 0) {$\x$};
}
% Z bei 0
\fill (0, 0) circle (2pt);
\node[above=10pt, font=\small] at (0, 0) {$Z$};
% P-Punkte links
\fill (-1, 0) circle (2pt)
    node[below=6pt, fill=white, inner sep=1pt, font=\small] {$P_{-1}$};
\fill (-3, 0) circle (2pt)
    node[below=6pt, fill=white, inner sep=1pt, font=\small] {$P_{-2}$};
% P-Punkte rechts
\fill (1, 0) circle (2pt)
    node[below=6pt, fill=white, inner sep=1pt, font=\small] {$P_1$};
\fill (3, 0) circle (2pt)
    node[below=6pt, fill=white, inner sep=1pt, font=\small] {$P_2$};
\fill (5, 0) circle (2pt)
    node[below=6pt, fill=white, inner sep=1pt, font=\small] {$P_3$};
\fill (7, 0) circle (2pt)
    node[below=6pt, fill=white, inner sep=1pt, font=\small] {$P_4$};
\fill (9, 0) circle (2pt)
    node[below=6pt, fill=white, inner sep=1pt, font=\small] {$P_5$};

\end{tikzpicture}
\end{document}

